%=====================================================================
\chapter{Measuring Classifier Performance}\label{ch:evaluation}
%=====================================================================
\locator{2}{Part II --- Linear Models and Generalization}

\begin{outcomes}
By the end of this chapter you will be able to:
\begin{tight}
  \item \textbf{Build} a confusion matrix and \textbf{compute} accuracy,
        precision, recall, specificity and $F_1$ from it.
  \item \textbf{Explain} why accuracy misleads on imbalanced classes, and
        \textbf{choose} a metric that matches the cost of each error.
  \item \textbf{Read} ROC and precision--recall curves, and \textbf{choose} a
        decision threshold from costs.
  \item \textbf{Compute} macro- and micro-averaged metrics for several classes.
  \item \textbf{Construct} a confidence interval for an error rate, and
        \textbf{test} whether two classifiers genuinely differ.
\end{tight}
\end{outcomes}

\begin{prereq}
\chref{math} (the base-rate calculation), \chref{logreg} (probabilities and
thresholds) and \chref{generalization} (cross-validation; the test set used
once).
\end{prereq}

\begin{keyterms}
confusion matrix $\bullet$ true/false positive/negative $\bullet$ accuracy
$\bullet$ precision $\bullet$ recall (sensitivity) $\bullet$ specificity
$\bullet$ $F_1$ score $\bullet$ class imbalance $\bullet$ ROC curve $\bullet$
AUC $\bullet$ precision--recall curve $\bullet$ operating threshold $\bullet$
macro and micro averaging $\bullet$ sample error $\bullet$ true error
$\bullet$ confidence interval $\bullet$ McNemar's test
\end{keyterms}

\section{Measuring Classifier Accuracy Is Not One Number}

The HEC outline calls this topic \emph{measuring classifier accuracy}, and the
first lesson is that accuracy is usually the wrong measure. A classifier makes
two kinds of mistake, and in almost every real problem one costs far more than
the other: a missed intrusion against a wasted analyst hour, an undetected
failing student against an unnecessary phone call. A single number that adds the
two together hides exactly the information needed to decide whether the
classifier is any use.

\section{The Confusion Matrix}

\begin{definitionbox}[Confusion Matrix]
For two classes, cross-tabulate the predictions against the truth:
\begin{tight}
  \item \term{true positive} (TP): positive, predicted positive;
  \item \term{false positive} (FP): negative, predicted positive --- a false alarm;
  \item \term{false negative} (FN): positive, predicted negative --- a miss;
  \item \term{true negative} (TN): negative, predicted negative.
\end{tight}
Every metric in this chapter is a ratio of these four counts.
\end{definitionbox}

\dsfig{%
\begin{tikzpicture}[font=\scriptsize,
  cell/.style={rounded corners=3pt, minimum width=2.9cm, minimum height=1.5cm,
               align=center, line width=0.9pt}]
\node[font=\scriptsize\bfseries, text=primary] at (1.55,2.0) {truly malicious};
\node[font=\scriptsize\bfseries, text=primary] at (4.65,2.0) {truly clean};
\node[font=\scriptsize\bfseries, text=primary, rotate=90] at (-0.35,0.8) {flagged};
\node[font=\scriptsize\bfseries, text=primary, rotate=90] at (-0.35,-0.8) {not flagged};
\node[cell, draw=greenhl, fill=greenhl!14] at (1.55,0.8)
     {{\color{greenhl!45!black}\bfseries TP = 225}\\[2pt]{\tiny caught}};
\node[cell, draw=accent, fill=accent!10] at (4.65,0.8)
     {{\color{accent}\bfseries FP = 75}\\[2pt]{\tiny false alarm}};
\node[cell, draw=accent, fill=accent!10] at (1.55,-0.8)
     {{\color{accent}\bfseries FN = 25}\\[2pt]{\tiny missed}};
\node[cell, draw=greenhl, fill=greenhl!14] at (4.65,-0.8)
     {{\color{greenhl!45!black}\bfseries TN = 4{,}675}\\[2pt]{\tiny left alone}};
\node[font=\tiny, text=gray!55!black] at (1.55,-1.8) {250 malicious};
\node[font=\tiny, text=gray!55!black] at (4.65,-1.8) {4{,}750 clean};
\node[font=\tiny, text=gray!55!black, anchor=west] at (6.2,0.8) {300 flagged};
\node[font=\tiny, text=gray!55!black, anchor=west] at (6.2,-0.8) {4{,}700 not};
% the metrics
\node[anchor=north west] at (7.9,2.2) {%
  \renewcommand{\arraystretch}{2.1}%
  \begin{tabular}{@{}l@{\quad}l@{\quad}l@{}}
  {\color{primary}\bfseries Accuracy} & $=\dfrac{\text{TP}+\text{TN}}{\text{all}}$ &
    {\tiny right overall} \\
  {\color{greenhl!45!black}\bfseries Precision} & $=\dfrac{\text{TP}}{\text{TP}+\text{FP}}$ &
    {\tiny of the flags, how many real} \\
  {\color{secondary}\bfseries Recall} & $=\dfrac{\text{TP}}{\text{TP}+\text{FN}}$ &
    {\tiny of the real ones, how many flagged} \\
  {\color{goldhl!50!black}\bfseries Specificity} & $=\dfrac{\text{TN}}{\text{TN}+\text{FP}}$ &
    {\tiny of the clean, how many left alone} \\
  {\color{purplehl}\bfseries $F_1$} & $=\dfrac{2\cdot\text{prec}\cdot\text{rec}}{\text{prec}+\text{rec}}$ &
    {\tiny harmonic mean of the two}
  \end{tabular}};
\end{tikzpicture}}%
{A malware scanner's confusion matrix on 5{,}000 files, and the five metrics
computed from it. Precision reads along the ``flagged'' row; recall reads down the
``truly malicious'' column.}{confusion}

\begin{worked}[One Scanner, Five Metrics]
The scanner of \dsref{confusion} examined 5{,}000 files, of which 250 were
malicious. It flagged 300, of which 225 were malicious.

\smallskip
\textbf{Cells:} TP $= 225$, FP $= 300 - 225 = 75$, FN $= 250 - 225 = 25$,
TN $= 5000 - 225 - 75 - 25 = 4675$.

\smallskip
\begin{tight}
  \item \textbf{Accuracy} $= (225 + 4675)/5000 = \mathbf{98.0\%}$
  \item \textbf{Precision} $= 225/300 = \mathbf{75.0\%}$
  \item \textbf{Recall} $= 225/250 = \mathbf{90.0\%}$
  \item \textbf{Specificity} $= 4675/4750 = \mathbf{98.4\%}$
  \item $\mathbf{F_1} = \dfrac{2(0.75)(0.90)}{0.75 + 0.90} = \dfrac{1.35}{1.65} = \mathbf{0.818}$
\end{tight}

\smallskip
\textbf{Against the baseline.} ``Flag nothing'' is right on the 4{,}750 clean
files: 95.0\% accuracy. The scanner's 98.0\% is three points better than
doing nothing, which sounds small; its recall of 90\% --- nine malicious files in
ten caught, at the cost of one false alarm for every three real detections --- is
what actually describes it.
\end{worked}

\begin{alertbox}[The Accuracy Paradox]
When one class is rare, accuracy is dominated by the common class. A detector
for an event that occurs in 1\% of cases can be 99\% accurate by never detecting
anything. Compute the majority-class baseline before looking at any accuracy
figure, and report precision and recall instead.
\end{alertbox}

\section{Precision, Recall and the Threshold}

A probabilistic classifier (\chref{logreg}) does not flag or clear anything by
itself; it outputs a score, and a \term{threshold} turns scores into decisions.
Lower the threshold and more cases are flagged: recall rises and precision falls.
Raise it and the reverse happens. Neither can be maximised without sacrificing
the other, so the threshold must be chosen from the costs of the two errors ---
which is a decision about the application, not about the algorithm.

\begin{worked}[Choosing an Operating Threshold From Costs]
For the malware scanner, a missed malicious file (FN) is estimated to cost 50
times as much as a false alarm (FP) --- an infection against ten minutes of an
analyst's time. The scanner's precision and recall at four thresholds, with 250
malicious files:

\smallskip
\begin{center}
\begin{tabular}{@{}rrrrrrr@{}}
\toprule
\textbf{Threshold} & \textbf{Precision} & \textbf{Recall} & \textbf{TP} & \textbf{FN} & \textbf{FP} & \textbf{Cost} $= 50\,\text{FN} + \text{FP}$ \\
\midrule
0.3 & 0.55 & 0.96 & 240 & 10 & 196 & \textbf{696} \\
0.5 & 0.75 & 0.90 & 225 & 25 & 75 & 1{,}325 \\
0.7 & 0.86 & 0.78 & 195 & 55 & 32 & 2{,}782 \\
0.9 & 0.95 & 0.52 & 130 & 120 & 7 & 6{,}007 \\
\bottomrule
\end{tabular}
\end{center}
\smallskip
(TP $=$ recall $\times$ 250; FP $=$ TP $\times (1 - \text{precision})/\text{precision}$.)

\smallskip
\textbf{The cheapest threshold is 0.3}, which almost halves the cost of the
default 0.5 by accepting 121 more false alarms to catch 15 more infections. With
a cost ratio of 1:1 the answer would be different. The threshold is where the
business decision enters the model.
\end{worked}

\section{Threshold-Free Summaries: ROC and PR Curves}

\begin{definitionbox}[ROC Curve and AUC]
The \term{ROC curve} plots the true positive rate (recall) against the false
positive rate ($1 -$ specificity) at every possible threshold. The \term{area
under the curve} (AUC) summarises it: the probability that the classifier scores
a randomly chosen positive above a randomly chosen negative. A random classifier
has AUC 0.5; a perfect one, 1.0.
\end{definitionbox}

\dsfig{%
\begin{tikzpicture}
\begin{groupplot}[
  group style={group size=2 by 1, horizontal sep=1.5cm},
  width=6.6cm, height=5.4cm, axis lines=left,
  tick label style={font=\tiny}, label style={font=\tiny},
  title style={font=\scriptsize\bfseries},
  xmin=0, xmax=1, ymin=0, ymax=1.03]
\nextgroupplot[title={ROC curves}, xlabel={false positive rate},
  ylabel={true positive rate (recall)},
  legend style={font=\tiny, draw=gray!40, at={(0.98,0.03)}, anchor=south east}]
% y = x^k has area 1/(1+k)
\addplot[name path=a, greenhl!70!black, line width=1.3pt, domain=0:1, samples=160] {x^(1/19)};
\addlegendentry{model A, AUC 0.95}
\addplot[goldhl, line width=1.2pt, domain=0:1, samples=160] {x^(1/3)};
\addlegendentry{model B, AUC 0.75}
\addplot[gray!55, dashed] coordinates {(0,0)(1,1)};
\addlegendentry{random, AUC 0.50}
\addplot[name path=z, draw=none, domain=0:1] {0};
\addplot[greenhl!8] fill between[of=a and z];
\node[font=\tiny, text=gray!55!black, anchor=north west] at (axis cs:0.03,0.99)
     {perfect corner};
\nextgroupplot[title={Precision--recall at 5\% positives}, xlabel={recall},
  ylabel={precision}]
\addplot[greenhl!70!black, line width=1.3pt, domain=0:1, samples=80] {0.95-0.55*x^3};
\addplot[goldhl, line width=1.2pt, domain=0:1, samples=80] {0.62-0.5*x^1.4};
\addplot[gray!55, dashed] coordinates {(0,0.05)(1,0.05)};
\node[font=\tiny, text=gray!55!black, anchor=south east] at (axis cs:0.98,0.06)
     {random: precision $=$ 5\%};
\node[font=\tiny, text=greenhl!45!black, anchor=south west] at (axis cs:0.05,0.95) {model A};
\node[font=\tiny, text=goldhl!45!black, anchor=south west] at (axis cs:0.05,0.63) {model B};
\end{groupplot}
\end{tikzpicture}}%
{Left: ROC curves sweep every threshold; the shaded area is model A's AUC. Right: the
same two models on precision and recall. With rare positives the PR plot shows how many
alerts are false, which the ROC plot hides.}{roc-pr}

\begin{alertbox}[Prefer Precision--Recall When Positives Are Rare]
The ROC false positive rate divides false alarms by the (huge) number of
negatives, so hundreds of false alarms can look like a rate of 0.01. A
precision--recall curve uses precision directly and does not flatter a
rare-event detector. For intrusion, fraud and failure detection, report the PR
curve or its area (average precision), and quote precision at the recall the
application needs.
\end{alertbox}

\section{More Than Two Classes}

With $k$ classes the confusion matrix is $k \times k$, and precision and recall
are computed for each class in turn, treating it as positive and the rest as
negative. They are then combined in one of two ways. \term{Macro averaging} takes
the plain mean of the per-class values, so every class counts equally.
\term{Micro averaging} pools the counts first, so every \emph{example} counts
equally --- and micro-averaged recall equals accuracy.

\begin{worked}[Macro Against Micro]
A traffic classifier is tested on 100 flows: 80 web, 15 e-mail, 5 streaming. It
labels 76 web, 9 e-mail and 1 streaming flow correctly.

\smallskip
\textbf{Per-class recall:} web $76/80 = 0.95$; e-mail $9/15 = 0.60$;
streaming $1/5 = 0.20$.

\textbf{Micro (accuracy):} $(76 + 9 + 1)/100 = \mathbf{0.86}$.

\textbf{Macro:} $(0.95 + 0.60 + 0.20)/3 = \mathbf{0.58}$.

\smallskip
The micro figure is carried by the web class; the macro figure reveals that the
classifier almost never recognises streaming. If the rare classes matter --- and
in security they usually are the ones that matter --- report the macro figure
and the per-class table.
\end{worked}

\section{How Sure Is the Number? Confidence Intervals for Error}

A test set is a sample, so the error measured on it is an estimate. Mitchell
distinguishes the \term{sample error} $\text{error}_S(h)$, the fraction of the $n$
test examples misclassified, from the \term{true error} $\text{error}_{\mathcal{D}}(h)$
over the whole population. Because each test example is an independent yes/no
trial, the number of mistakes is binomial, and for $n \ge 30$ an approximate 95\%
confidence interval for the true error is
\[
\text{error}_S(h) \;\pm\; 1.96\,\sqrt{\frac{\text{error}_S(h)\,\big(1 - \text{error}_S(h)\big)}{n}}
\]

\begin{worked}[How Wide Is 12\% Error?]
A classifier misclassifies 24 of $n = 200$ test examples: $\text{error}_S = 0.12$.
\[
\sqrt{\frac{0.12 \times 0.88}{200}} = \sqrt{0.000528} = 0.0230,
\qquad 1.96 \times 0.0230 = 0.045
\]
\textbf{95\% interval:} $0.12 \pm 0.045$, that is \textbf{7.5\% to 16.5\%}.

\smallskip
With $n = 2000$ and the same 12\% error, the half-width shrinks to
$1.96\sqrt{0.12 \times 0.88/2000} = 0.014$: \textbf{10.6\% to 13.4\%}. Ten times
the data gives an interval about $\sqrt{10} \approx 3.2$ times narrower. On 200
examples, a competitor at 10\% error is not demonstrably better.
\end{worked}

\dsfig{%
\begin{tikzpicture}
\begin{axis}[width=9.4cm, height=4.8cm, axis lines=left, xmode=log,
  xlabel={test set size $n$}, ylabel={95\% interval for error},
  tick label style={font=\tiny}, label style={font=\scriptsize},
  xmin=40, xmax=12000, ymin=0, ymax=0.25, log basis x=10,
  yticklabel style={/pgf/number format/fixed}]
\addplot[name path=up, secondary, line width=1.1pt, domain=40:12000, samples=80]
  {0.12+1.96*sqrt(0.12*0.88/x)};
\addplot[name path=lo, secondary, line width=1.1pt, domain=40:12000, samples=80]
  {0.12-1.96*sqrt(0.12*0.88/x)};
\addplot[secondary!12] fill between[of=up and lo];
\addplot[primary, dashed, line width=0.9pt] coordinates {(40,0.12)(12000,0.12)};
\addplot[only marks, mark=*, mark size=1.8pt, accent] coordinates {(200,0.075)(200,0.165)(2000,0.106)(2000,0.134)};
\node[font=\tiny, accent, anchor=north west] at (axis cs:215,0.068) {$n = 200$: 7.5--16.5\%};
\node[font=\tiny, accent, anchor=south east] at (axis cs:1900,0.168) {$n = 2000$: 10.6--13.4\%};
\draw[accent, line width=0.4pt] (axis cs:1900,0.17) -- (axis cs:2000,0.136);
\node[font=\tiny, primary, anchor=south west] at (axis cs:45,0.121) {measured error 12\%};
\end{axis}
\end{tikzpicture}}%
{How the uncertainty in a measured 12\% error rate shrinks with test-set size. The width
falls as $1/\sqrt{n}$: four times the data halves it.}{error-ci}

\section{Is One Classifier Really Better? McNemar's Test}

Two classifiers are tested on the same examples and one scores a little higher.
Is it better, or luckier? Only the examples on which they \emph{disagree} carry
information. Let $b$ be the number that classifier A gets wrong and B gets
right, and $c$ the number A gets right and B wrong. If the two are equally good,
$b$ and $c$ should be similar. \term{McNemar's test} statistic
\[
\chi^2 = \frac{(|b - c| - 1)^2}{b + c}
\]
is compared with 3.84, the 95\% point of the chi-squared distribution with one
degree of freedom; a larger value means the difference is unlikely to be chance.

\begin{worked}[McNemar on 500 Test Emails]
A spam classifier A and a new classifier B are run on the same 500 e-mails. They
disagree on 40: B is right and A wrong on $b = 12$; A is right and B wrong on
$c = 28$.
\[
\chi^2 = \frac{(|12 - 28| - 1)^2}{12 + 28} = \frac{15^2}{40} = 5.63 \;>\; 3.84
\]
\textbf{A is significantly better} at the 5\% level ($p \approx 0.018$): it is right on
16 more e-mails, 3.2 points of accuracy. The 460 e-mails on which they agree play
no part --- which is why the test is sensitive even when the accuracies look
close. When the comparison is across cross-validation folds instead of one test
set, the equivalent tool is a paired $t$-test on the per-fold scores.
\end{worked}

\section{Handling Class Imbalance}

Imbalance is a problem for training as well as evaluation: a learner minimising
average loss can do very well by ignoring the rare class. Three remedies, in
order of preference:
\begin{tightnum}
  \item \textbf{Move the threshold.} Train normally, then choose the threshold from
        costs, as above. It changes nothing about the model and is easy to
        revisit.
  \item \textbf{Weight the classes.} Make each error on the rare class cost more in
        the loss (\texttt{class\_weight="balanced"} in scikit-learn).
  \item \textbf{Resample the training data.} Undersample the common class, or
        oversample the rare one (by duplication, or by synthesising new examples
        with SMOTE). Resample only the training folds --- never the validation or
        test data, whose class balance must stay real.
\end{tightnum}

\begin{pitfall}
\begin{tight}
  \item \textbf{Reporting accuracy on imbalanced data} without the majority
        baseline beside it.
  \item \textbf{Quoting precision without recall,} or either without the
        threshold at which it was measured.
  \item \textbf{Using ROC-AUC alone for rare events.} Report the PR curve.
  \item \textbf{Declaring a winner from a one-point difference} on a few hundred
        test examples. Compute the interval or run McNemar's test.
  \item \textbf{Resampling before splitting,} which puts copies of the same rare
        example into both training and test sets.
\end{tight}
\end{pitfall}

%---------------------------------------------------------------------
\begin{fourlenses}{Choosing the Metric}
\lensLA{Recall at a fixed outreach budget: the advisers can call 50 students, so
report how many of the eventual failures are among the top 50 scores. Precision
matters too --- if most calls reach students who were fine, the advisers stop
trusting the list.}

\lensPM{Over- and under-predicting a project overrun cost roughly the same, so
$F_1$ or balanced accuracy is reasonable. With 40 projects in the test set, the
confidence interval on any rate is about $\pm 15$ points; say so, rather than
reporting three significant figures.}

\lensIOT{Failures are rare and a miss can stop a production line, so optimise
recall and then cap false alarms at what the maintenance team can inspect. Use
the PR curve; ROC-AUC will look excellent on a detector the technicians
quickly learn to ignore.}

\lensSEC{Report precision at a fixed recall --- ``62\% precision at 90\%
recall'' --- and alerts per analyst per hour. Choose the threshold from the cost
ratio, as in the malware example; and compare a new detector with the old on the
same traffic with McNemar's test, not on two different weeks.}
\end{fourlenses}

%---------------------------------------------------------------------
\begin{lab}[Evaluate an Imbalanced Classifier Properly]
\begin{lstlisting}[language=Python]
import numpy as np
from scipy.stats import chi2
from sklearn.datasets import make_classification
from sklearn.model_selection import train_test_split
from sklearn.linear_model import LogisticRegression
from sklearn.tree import DecisionTreeClassifier
from sklearn.metrics import (confusion_matrix, precision_score, recall_score,
                             f1_score, roc_auc_score, average_precision_score)

# --- an imbalanced problem: 5% positives ---------------------------------
X, y = make_classification(n_samples=6000, n_features=12, n_informative=5,
                           weights=[0.95], flip_y=0.01, random_state=1)
X_tr, X_te, y_tr, y_te = train_test_split(X, y, test_size=0.4,
                                          stratify=y, random_state=1)
print("baseline accuracy (always 0):", round(1 - y_te.mean(), 3))
lr = LogisticRegression(max_iter=1000).fit(X_tr, y_tr)
p = lr.predict_proba(X_te)[:, 1]

# --- 1. the four cells and the metrics built from them -------------------
pred = (p > 0.5).astype(int)
tn, fp, fn, tp = confusion_matrix(y_te, pred).ravel()
print(f"TP={tp} FP={fp} FN={fn} TN={tn}")
print(f"accuracy {np.mean(pred == y_te):.3f}  "
      f"precision {precision_score(y_te, pred):.3f}  "
      f"recall {recall_score(y_te, pred):.3f}  F1 {f1_score(y_te, pred):.3f}")
print(f"ROC-AUC {roc_auc_score(y_te, p):.3f}   "
      f"PR-AUC {average_precision_score(y_te, p):.3f}")

# --- 2. the threshold is a choice: sweep it ------------------------------
for t in (0.1, 0.3, 0.5, 0.7):
    pr = (p > t).astype(int)
    print(f"threshold {t}: precision {precision_score(y_te, pr):.2f}  "
          f"recall {recall_score(y_te, pr):.2f}")

# --- 3. a 95% confidence interval for the error rate ---------------------
n, e = len(y_te), np.mean(pred != y_te)
half = 1.96 * np.sqrt(e * (1 - e) / n)
print(f"error {e:.3f}, 95% interval [{e - half:.3f}, {e + half:.3f}]")

# --- 4. is a decision tree really different? McNemar's test --------------
tree = DecisionTreeClassifier(max_depth=5, random_state=0).fit(X_tr, y_tr)
a_right = pred == y_te
b_right = tree.predict(X_te) == y_te
b = np.sum(~a_right & b_right)          # logistic wrong, tree right
c = np.sum(a_right & ~b_right)          # logistic right, tree wrong
stat = (abs(b - c) - 1) ** 2 / (b + c)
print(f"disagreements b={b}, c={c}; McNemar chi2 = {stat:.2f}, "
      f"p = {chi2.sf(stat, 1):.3f}")
\end{lstlisting}
The model is 96\% accurate against a 95\% baseline, and at the default threshold it
finds only a third of the positives. Add \texttt{class\_weight="balanced"} to
\texttt{LogisticRegression} and compare the recall, precision and PR-AUC.
\end{lab}

%---------------------------------------------------------------------
\begin{checkpoint}
\begin{tightnum}
  \item Of 400 flagged logins, 100 were attacks; 20 attacks were not flagged.
        Compute precision and recall.
  \item A classifier's ROC-AUC is 0.97 on data with 0.1\% positives. What should
        you ask for before believing it is useful?
  \item A model makes 15 errors on 100 test examples. Give an approximate 95\%
        interval for its true error.
\end{tightnum}
\end{checkpoint}

%---------------------------------------------------------------------
\begin{chaptersummary}
\begin{tight}
  \item Every classification metric is a ratio of the four confusion-matrix
        cells; always compute the majority baseline first.
  \item Precision asks ``when I flag, am I right?''; recall asks ``of the real
        cases, how many did I flag?'' The threshold trades one for the other,
        and should be chosen from the costs of the two errors.
  \item ROC-AUC summarises ranking quality across thresholds but flatters
        rare-event detectors; prefer precision--recall curves there.
  \item Macro averaging treats classes equally and exposes weak rare classes;
        micro averaging equals accuracy.
  \item A measured error is an estimate: its 95\% interval is
        $e \pm 1.96\sqrt{e(1-e)/n}$. Compare two classifiers on the same test set
        with McNemar's test.
  \item Handle imbalance by moving the threshold, weighting classes, or
        resampling the training folds only.
\end{tight}
\end{chaptersummary}

\begin{reviewq}
  \item \textbf{(MCQ)} Recall is: (a)~TP/(TP+FP) (b)~TP/(TP+FN) (c)~TN/(TN+FP)
        (d)~(TP+TN)/all.
  \item \textbf{(MCQ)} Lowering the decision threshold usually: (a)~raises
        precision and recall (b)~raises recall, lowers precision (c)~raises
        precision, lowers recall (d)~changes neither.
  \item \textbf{(MCQ)} Micro-averaged recall over all classes equals:
        (a)~macro recall (b)~accuracy (c)~precision (d)~$F_1$.
  \item \textbf{(Short)} Explain the accuracy paradox with a numeric example of
        your own.
  \item \textbf{(Short)} Why do only the disagreements matter in McNemar's test?
  \item \textbf{(Short)} Give a situation where you would choose a threshold that
        makes precision 0.2, and justify it.
  \item \textbf{(Applied)} A model tested on 10{,}000 connections, 200 of them
        attacks, flags 500, of which 160 are attacks. Build the confusion matrix
        and compute accuracy, precision, recall, specificity and $F_1$. Which two
        numbers would you put in a report to the security manager?
  \item \textbf{(Applied)} Classifier A has 18\% error and classifier B 15\% on the
        same 300 test examples. Compute a 95\% interval for each. On the
        disagreements, A is wrong and B right 25 times, and the reverse 16
        times. Is B significantly better?
\end{reviewq}
